\documentclass[12pt]{article}
\usepackage{amsmath,amssymb,amsfonts}
\usepackage[margin=1in]{geometry}

\begin{document}

\section*{Chapter 7, Problem 28}

\textbf{Problem:} Using the Fourier transform method (in four dimensions) obtain the Green's functions for the wave equation. Show that one obtains $G_A$ of Eq.\ (7.163) or $G_R$ (of the previous problem) depending on how one chooses to avoid a singularity which arises in the Fourier transform method (see a similar discussion in Section 7.4, Eq.\ (7.43) et seq.).

\bigskip
\textbf{Solution:}

The wave equation Green's function satisfies:
\[
\Bigl(\nabla^2 - \frac{1}{v^2}\frac{\partial^2}{\partial t^2}\Bigr)G = \delta^3(\mathbf{r}-\mathbf{r}')\delta(t-t').
\]

Take the four-dimensional Fourier transform. Let $\mathbf{R} = \mathbf{r}-\mathbf{r}'$ and $\tau = t-t'$:
\[
\tilde{G}(\mathbf{k},\omega) = \int G(\mathbf{R},\tau)\,e^{-i\mathbf{k}\cdot\mathbf{R}+i\omega\tau}\,d^3R\,d\tau.
\]

The equation becomes:
\[
(-k^2 + \omega^2/v^2)\tilde{G} = 1,
\]
so:
\[
\tilde{G}(\mathbf{k},\omega) = \frac{1}{\omega^2/v^2 - k^2} = \frac{v^2}{\omega^2 - v^2k^2}.
\]

Inverting:
\[
G(\mathbf{R},\tau) = \int\frac{d^3k}{(2\pi)^3}\int\frac{d\omega}{2\pi}\,\frac{v^2\,e^{i\mathbf{k}\cdot\mathbf{R}-i\omega\tau}}{\omega^2 - v^2k^2}.
\]

The $\omega$-integral has poles at $\omega = \pm vk$ on the real axis. The choice of contour determines the Green's function:

\textbf{Retarded:} Push both poles below the real axis ($\omega \to \omega + i\epsilon$, i.e., $\omega = \pm vk - i\epsilon$). For $\tau > 0$, close in the lower half-plane, picking up both poles:
\[
\int\frac{e^{-i\omega\tau}}{\omega^2 - v^2k^2}\,\frac{d\omega}{2\pi} = \int\frac{e^{-i\omega\tau}}{(\omega-vk)(\omega+vk)}\,\frac{d\omega}{2\pi} = -\frac{\sin(vk\tau)}{vk}\,\theta(\tau).
\]

For $\tau < 0$, close in the upper half-plane --- no poles, giving zero.

\textbf{Advanced:} Push both poles above the real axis ($\omega = \pm vk + i\epsilon$). This gives the same result but with $\theta(-\tau)$ instead.

Continuing with the retarded case:
\[
G_R(\mathbf{R},\tau) = -\theta(\tau)\int\frac{d^3k}{(2\pi)^3}\,\frac{v\sin(vk\tau)}{k}\,e^{i\mathbf{k}\cdot\mathbf{R}}.
\]

Performing the angular integration ($\mathbf{k}\cdot\mathbf{R} = kR\cos\alpha$):
\[
\int\frac{d^3k}{(2\pi)^3}\frac{\sin(vk\tau)}{k}\,e^{i\mathbf{k}\cdot\mathbf{R}} = \frac{1}{2\pi^2 R}\int_0^{\infty}\sin(vk\tau)\sin(kR)\,dk.
\]

Using $\sin A\sin B = \frac{1}{2}[\cos(A-B)-\cos(A+B)]$ and $\int_0^\infty\cos(ax)\,dk = \pi\delta(a)$:
\[
= \frac{1}{4\pi R}\bigl[\delta(v\tau - R) - \delta(v\tau + R)\bigr].
\]

Since $\tau > 0$ and $R > 0$, the second delta never contributes, so:
\[
G_R = -\theta(\tau)\cdot\frac{v}{4\pi R}\,\delta(v\tau - R) = -\frac{\delta(\tau - R/v)}{4\pi R}.
\]

This is the retarded Green's function. Similarly, the advanced prescription gives:
\[
G_A = -\frac{\delta(\tau + R/v)}{4\pi R},
\]
confirming Eq.\ (7.163). The ambiguity in the Fourier method corresponds exactly to the physical choice of boundary conditions (outgoing vs.\ incoming waves). $\blacksquare$

\end{document}
